\section{The ordinary endpoint trace of the actual Weil eigenfunctions}
\label{sec:spectral-endpoint}

The boundary functional used in the finite spectral triples must
be distinguished from the ordinary trace of an eigenfunction of
the closed interval operator.  The distinction is substantive
already for a fixed interval.

\begin{theorem}[AC-S36: bounded eigenfunctions and zero ordinary boundary values]
\label{AC-S36}
Fix \(L>0\), put \(I=(-L,L)\), and let \(A_L\) be the actual
closed zeta Weil operator of AC-S03, without assuming \(A\).
Every eigenfunction \(u\) of \(A_L\), extended by zero outside
\(I\), has a representative in \(C(\mathbb R)\).  For every
\(0<\tau<1/2\),
\[
 |u(x)|\leq C_{u,L,\tau}
          \left(\log\frac1{d(x)}\right)^{-\tau},
 \qquad d(x)=\operatorname{dist}(x,\partial I)\longrightarrow0,
 \quad x\in I.                                      \tag{EP.1}
\]
In particular,
\[
                         u(-L)=u(L)=0.              \tag{EP.2}
\]
More generally, the heat semigroup has the fixed-window
regularization
\[
              e^{-tA_L}:L^2(I)\longrightarrow L^\infty(I)
              \quad\hbox{boundedly whenever }t>1/2. \tag{EP.3}
\]
This is a domination statement, not positivity preservation
of the Weil semigroup; the latter fails in the range specified
by AC-S28.
\end{theorem}
\begin{proof}
\proofstep{1. The singular part and a finite convolution majorant.}
Use zero extension throughout and set
\[
 \mu(h)=\frac{e^{-|h|/2}}{1-e^{-2|h|}},\qquad
 a_0=-\log\pi+\psi(1/4).
\]
Let \(H_{\mu,I}\) be the nonnegative operator associated with
\[
 \mathfrak h_{\mu,I}(f)=\frac12
       \iint_{\mathbb R^2}\mu(x-y)|f(x)-f(y)|^2\,dx\,dy.
                                                        \tag{EP.4}
\]
The full-domain identity (SB.26) gives
\[
                      A_L=H_{\mu,I}+a_0+R_L,          \tag{EP.5}
\]
where, writing \(T_hf(x)=\widetilde f(x-h)\) on \(I\),
\[
 \begin{split}
 R_Lf(x)={}&-\sum_{2\leq n\leq e^{2L}}
       \frac{\Lambda(n)}{\sqrt n}
           (T_{\log n}f(x)+T_{-\log n}f(x))\\
 &+\int_I 2\cosh((x-y)/2)f(y)\,dy .
 \end{split}                                           \tag{EP.6}
\]
Here \(\Lambda(n)\) is the von Mangoldt function and a tilde
denotes zero extension.  The operator \(R_L\) is self-adjoint
and bounded on \(L^2(I)\), and bounded on every \(L^p(I)\),
\(1\leq p\leq\infty\).  Define the finite positive symmetric
measure
\[
 v_L=\sum_{2\leq n\leq e^{2L}}
       \frac{\Lambda(n)}{\sqrt n}
          (\delta_{\log n}+\delta_{-\log n})
       +1_{[-2L,2L]}(h)\,2\cosh(h/2)\,dh,
 \qquad M_L=v_L(\mathbb R).
                                                        \tag{EP.7}
\]
Directly from (EP.6),
\[
                    |R_Lf|\leq 1_I(v_L*|\widetilde f|).
                                                        \tag{EP.8}
\]
No sign assumption on \(f\), on the lowest eigenvalue, or
on the Weil form is involved.

\proofstep{2. Domination of the restricted jump semigroup.}
The whole-line operator \(H_\mu\) has Fourier symbol
\[
 \Psi_\mu(\xi)=\int_{\mathbb R}
            \mu(h)(1-\cos(\xi h))\,dh
       =a(\xi)-a_0,\quad
 a(\xi)=-\log\pi+
        \operatorname{Re}\psi(1/4+i\xi/2).
                                                        \tag{EP.9}
\]
The integral is finite because
\(\int\min(1,h^2)\mu(h)\,dh<\infty\).
Its semigroup \(S_\mu(t)=e^{-tH_\mu}\) is positive.
The following elementary approximation also verifies
the needed comparison with the restricted operator.

For \(\varepsilon>0\), put
\(\mu_\varepsilon=1_{|h|>\varepsilon}\mu\) and
\(m_\varepsilon=\int\mu_\varepsilon\).
Let \(P\) mean restriction followed by zero extension.
The restricted finite-jump operator is
\[
 H_{\varepsilon,I}
       =m_\varepsilon\,1-P(\mu_\varepsilon*)P.
\]
Its exponential has the norm-convergent series
\[
 e^{-tH_{\varepsilon,I}}
   =e^{-tm_\varepsilon}
        \sum_{k\geq0}\frac{t^k}{k!}
                    \bigl(P(\mu_\varepsilon*)P\bigr)^k.
\]
Every convolution kernel in this series is nonnegative.
Dropping intermediate restrictions can only increase
its value on a nonnegative function.  Thus
\[
 0\leq e^{-tH_{\varepsilon,I}}f
           \leq 1_I e^{-tH_{\mu_\varepsilon}}\widetilde f
                    \qquad(f\geq0).
\]
The nonnegative restricted forms increase to (EP.4)
as \(\varepsilon\downarrow0\).  Here is the variational
justification of the strong resolvent convergence being
used.  For fixed \(\sigma>0\) and \(f\in L^2(I)\), their
resolvent vectors \(u_\varepsilon\) minimize
\[
 J_\varepsilon(v)=\mathfrak h_{\varepsilon,I}(v)
              +\sigma\|v\|_2^2-2\operatorname{Re}\langle f,v\rangle .
\]
The resolvent equation bounds \(u_\varepsilon\) and its
energy uniformly.  If \(u\) is a weak subsequential limit,
then, for every fixed \(\delta>0\), weak lower
semicontinuity of the bounded nonnegative form
\(\mathfrak h_{\delta,I}\) gives
\[
 \mathfrak h_{\delta,I}(u)
        \leq\liminf_{\varepsilon\downarrow0}
                        \mathfrak h_{\varepsilon,I}(u_\varepsilon).
\]
Take the supremum over \(\delta\), and compare
\(J_\varepsilon(u_\varepsilon)\leq J_\varepsilon(v)\)
for each fixed vector in the limit form domain.
It follows that \(u\) is the unique minimizer of the
limit functional and the minimum values converge.
Since
\[
 \sigma\|u-u_\varepsilon\|_2^2
       \leq J_\varepsilon(u)-J_\varepsilon(u_\varepsilon)
                    \longrightarrow0,
\]
the resolvents converge strongly, without extraction.
The scalar functions \((1+tx/m)^{-m}\) converge to
\(e^{-tx}\) uniformly for \(x\geq0\).
Applying this Euler approximation to the nonnegative
operators, first fixing \(m\), proves strong convergence
of their semigroups as well.
On the whole line the corresponding
strong convergence follows directly from Plancherel and
dominated convergence of \(e^{-t\Psi_{\mu_\varepsilon}}\).
Taking an almost-everywhere convergent subsequence in the
last inequality therefore gives
\[
 0\leq e^{-tH_{\mu,I}}f
             \leq1_I S_\mu(t)\widetilde f
                    \qquad(f\geq0).                   \tag{EP.10}
\]
The same approximation proves positivity of \(S_\mu(t)\);
no independent assumption about a heat kernel is needed.

\proofstep{3. The bounded perturbation gives eventual \(L^2\)-to-\(L^\infty\)
regularization.}
Apply the iterated Duhamel expansion to the bounded
perturbation \(R_L\) of \(H_{\mu,I}\).  It converges in
operator norm on every bounded time interval.
By (EP.8), (EP.10), and positivity, the absolute value
of its term with \(k\) perturbations is bounded by
\[
 \frac{t^k}{k!}\,1_I
                 S_\mu(t)(v_L^{*k}*|\widetilde f|).
\]
Indeed each intermediate restriction may be removed;
the remaining whole-line convolution operators commute,
and the time simplex has volume \(t^k/k!\).
Summing these positive bounds gives
\[
 |e^{-tA_L}f|
     \leq e^{-ta_0}\,1_I S_\mu(t)
                  \left(\sum_{k\geq0}
                         \frac{t^k}{k!}v_L^{*k}
                         *|\widetilde f|\right).
                                                        \tag{EP.11}
\]
The positive measure in parentheses has total mass
\(e^{tM_L}\).  Its Fourier transform is
\(\exp(t\widehat v_L(\xi))\).

The digamma asymptotic in (EP.9) gives
\(\Psi_\mu(\xi)=\log(2+|\xi|)+O(1)\).
Hence \(e^{-t\Psi_\mu}\in L^2(\mathbb R)\) for \(t>1/2\).
Fourier inversion and Cauchy--Schwarz, with the
unnormalized Fourier convention of this paper, yield
\[
 \|e^{-tA_L}f\|_\infty
 \leq \frac{e^{t(M_L-a_0)}}{\sqrt{2\pi}}\,
          \|e^{-t\Psi_\mu}\|_{L^2(\mathbb R)}\|f\|_2,
                  \qquad t>1/2.                       \tag{EP.12}
\]
This proves (EP.3).  If \(A_Lu=\alpha u\), then
\(u=e^{t\alpha}e^{-tA_L}u\), so \(u\in L^\infty(I)\).
The constants here may depend on the fixed interval;
no uniform growing-window bound is being asserted.

\proofstep{4. Exact reduction to a logarithmic Dirichlet problem.}
Let \(L_\Delta\) denote the logarithmic Laplacian in
one dimension, with whole-line Fourier symbol
\(2\log|\xi|\) and integral normalization
\[
 \begin{split}
 L_\Delta f(x)={}&
   \int_{|h|<1}\frac{f(x)-f(x-h)}{|h|}\,dh
       -\int_{|h|\geq1}\frac{f(x-h)}{|h|}\,dh
       +\rho_1f(x),\\
 &\rho_1=-2\gamma_{\rm E}.
 \end{split}                                           \tag{EP.13}
\]
The integrals are understood in the weak form for the
zero extension, not assumed pointwise for an arbitrary
form vector.
The function
\[
 k(h)=\mu(h)-\frac{1_{0<|h|<1}}{2|h|}
\]
belongs to \(L^1(\mathbb R)\):
it is bounded near zero, and \(\mu\) decreases
exponentially at infinity.  Put
\[
 b_L(h)=\frac{1_{1\leq|h|\leq2L}}{2|h|},
 \qquad \kappa=\int_{\mathbb R}k(h)\,dh,
\]
with \(b_L=0\) if \(2L<1\).
On zero-extended functions supported in \(I\), direct
subtraction of the weak kernels gives
\[
 H_{\mu,I}
   =\tfrac12 L_{\Delta,I}
        +(\kappa-\rho_1/2)\,1
           -P(k*)P+P(b_L*)P.                          \tag{EP.14}
\]
For example, the final term restores the long-range
negative integral in (EP.13); shifts longer than \(2L\)
have zero compression.  Thus the signs and the factor
one half in (EP.14) are fixed by the displayed kernels.

Both forms in (EP.14) have the same domain:
their near-diagonal kernels differ by the integrable
function \(k\), and every remaining compressed term is
bounded on \(L^2(I)\).  Equivalently, this is the space
\(\mathcal H(I)\) in Chen--Weth, defined immediately
before their Theorem 1.4 by the zero-extension
quadratic integral with kernel
\(1_{|h|<1}/|h|\).
Combining (EP.5) and (EP.14), write
\[
                         A_L=\tfrac12L_{\Delta,I}+B_L,
                                                        \tag{EP.15}
\]
where \(B_L\) is bounded on both \(L^2(I)\) and
\(L^\infty(I)\).
For an eigenfunction \(u\), therefore,
\[
       L_\Delta u=2(\alpha u-B_Lu)\in L^\infty(I)
       \quad\hbox{weakly in }I,\qquad u=0\text{ off }I.
                                                        \tag{EP.16}
\]
All hypotheses on the solution and forcing have now
been proved, including boundedness of \(u\).

\proofstep{5. Boundary regularity and the source normalization.}
We use H.~Chen and T.~Weth,
\emph{The Dirichlet Problem for the Logarithmic Laplacian},
arXiv:1710.03416v6, Theorem 1.11, printed p.~5,
restated as Theorem 5.4, printed pp.~28--30,
\url{https://arxiv.org/pdf/1710.03416v6}.
That theorem gives interior continuity and (EP.1)
for a bounded weak solution with bounded forcing
on a Lipschitz domain with a uniform exterior sphere.
The interval \(I\) satisfies these geometric hypotheses.
It applies to (EP.16), separately to real and imaginary
parts if needed.
Interior continuity and (EP.1) together give the
continuous zero extension and (EP.2).

For the source multiplicative coordinate \(r=e^x\),
the continuous representative
\(\xi(r)=u(\log r)\) consequently satisfies
\[
                     \xi(\lambda^{-1})=\xi(\lambda)=0,
                     \qquad \lambda=e^L.              \tag{EP.17}
\]
This conclusion holds for every eigenfunction and
requires neither RH nor simplicity of the eigenspace.
\end{proof}

\begin{remark}[The exact consequence for the proposed determinant limit]
S25, the first item immediately before equation
\texttt{ku} in Section 7, describes a fixed-window limit
using a lowest eigenfunction normalized by
\(\xi_\lambda(\lambda)=1\).
Equation (EP.17) shows that this cannot be the ordinary
continuous endpoint trace of the actual closed Weil
operator.  Changing the value of an \(L^2\) representative
at a single point would not define such a normalization:
it changes neither the eigenvector nor its transform.
An alternative boundary trace would have to be defined
and proved to exist.

The finite-dimensional boundary functionals in S25,
\texttt{dirichletdef}, \texttt{dirichlet1}, and
\texttt{finmain}, remain well defined, and their
nonvanishing under the stated even-simple hypothesis
is unaffected.  AC-S06, with the semibounded extension
explained after AC-S28, gives the valid fixed-window
convergence after \(L^2\) normalization whenever the
closed lowest line is simple and even, including the
proved range of AC-S28.
This theorem does not assert that the original finite
boundary-normalized determinants diverge: convergence
in form norm or \(L^2\) does not by itself pass their
varying boundary functionals to the limit.
It also does not settle the growing-window convergence
of the lowest eigenline to \(K\).
\end{remark}
